% Supplement to source_and_masullo.tex. Original proofs of the stated deductions;
% no claim of historical novelty. Background inputs are explicitly identified.
\section{The actual Hilbert-space sheets and a numerical completion over \texorpdfstring{$\mathbb Q$}{Q}}
\label{sec:masullo-numeric}

\subsection{What cote 1 actually says}

The supplied transcription of cote~1, scan~079, gives Corollaries~4 and~5:
coefficientwise domination of the Fredholm determinant of $T$ by that of
$|T|$, and hence $|\det(1+zT)|\leq\det(1+|z||T|)$.  The following
eigenvalue/product theorem and its trace corollary are crossed out, as is the
draft proof continuing on scan~080.  Scans~111--112 state and attempt to prove
the genus-zero lemma used by that draft.  These supplied sheets contain no
Riemann-$\xi$ function, arithmetic spectral realization, or unfinished question
asserting that such a realization should exist.  Thus BR-TR-01 and BR-TR-02
are arithmetic applications of this analytic framework.  The next two
findings also give complete proofs of the actual analytic assertions on the
supplied sheets, including the crossed-out product assertion.  ``Crossed
out'' is a description of this source, not a claim that the classical theorem
is false or remains open.

\paragraph{Finding F-SM-11 (the proposition on scans 111--112, corrected).}
Let $f(z)=\sum_{n\geq0}a_nz^n\not\equiv0$ and
$h(z)=\sum_{n\geq0}b_nz^n\not\equiv0$ be entire, with
$|a_n|\leq|b_n|$ for all $n$.  Suppose that $h$ has genus zero, and that the
nonzero zeros $z_j$ of $f$, repeated with multiplicity, satisfy
$\sum_j|z_j|^{-1}<\infty$.  Then $f$ has genus zero.  Here genus zero means
a constant times a power of $z$ times a normally convergent product of
factors $1-z/z_j$, with no nonconstant exponential factor.

\begin{proof}
Write $h(z)=c z^r\prod_j(1-z/w_j)$, where
$\sum_j|w_j|^{-1}<\infty$.  The function
$H(z)=|c|z^r\prod_j(1+z/|w_j|)$ has nonnegative coefficients that majorize
the absolute values of the coefficients of $h$.  To see this, first use the
finite products, then pass to their locally uniform limits; Cauchy's
coefficient formula preserves each coefficient inequality.  Therefore
$|f(z)|\leq H(|z|)$.

Every genus-zero product $P$ satisfies, for each $\varepsilon>0$,
$|P(z)|\leq C_\varepsilon e^{\varepsilon|z|}$.  Indeed, split its factors
after an index $N$ for which the sum of the reciprocal moduli in the tail is
less than $\varepsilon/2$.  The tail is bounded by
$\exp((\varepsilon/2)|z|)$, using $\log(1+x)\leq x$; the finite product
and the power of $z$ are bounded by a constant times
$\exp((\varepsilon/2)|z|)$.  Applying this to $H$ proves the same bound
for $f$.  Hadamard factorization for an entire function of order at most one
now gives
\[
 f(z)=e^{A+bz}P(z),\qquad P(z)=z^m\prod_j(1-z/z_j).
\]
More explicitly, the genus-one factors in the usual Hadamard formula may
be replaced by genus-zero factors, because $\sum_j1/z_j$ converges
absolutely; the resulting linear exponential is absorbed into $b$.

Suppose $b\ne0$.  Replace $z$ by $\overline b z/|b|$ to make $b>0$;
the preceding radial bounds are unchanged.  Fix
$\pi/3<\theta<\pi/2$, and put $\delta=b\cos\theta/2>0$.  On the sector
$|\arg z|\leq\theta$, the bound for $f$ with $\varepsilon=\delta$ gives
\[
 |P(z)|\leq C e^{-\delta|z|}.
\]
Globally the genus-zero bound for $P$ gives
$|P(z)|\leq C'e^{\delta|z|/3}$.  Set $\omega=e^{2\pi i/3}$.
For every nonzero $z$, at least one of $z,\omega z,\omega^2z$ belongs
to this sector, because its half-width exceeds $\pi/3$.  Hence the entire
function $G(z)=P(z)P(\omega z)P(\omega^2z)$ satisfies
\[
 |G(z)|\leq C(C')^2e^{-\delta|z|/3}.
\]
Liouville's theorem makes $G$ constant, and the displayed bound makes that
constant zero.  A finite product of nonzero entire functions cannot vanish
identically.  This contradicts $P\not\equiv0$, so $b=0$.
\end{proof}

Three transcription-level points matter.  Scan~112 literally specifies
$\pi/3<\theta<\pi$ while subsequently requiring
$\Re z\geq c|z|$ throughout the sector, for $c>0$.
The latter assertion requires $\theta<\pi/2$; the proof above makes the
compatible choice $\pi/3<\theta<\pi/2$.  Scan~111 writes a rotation by
$b/|b|$; with the substitution $f(\Theta z)$ the correcting rotation is
$\Theta=\overline b/|b|$.  Finally, coefficient domination gives
$|f(z)|\leq H(|z|)$, not generally $|f(z)|\leq|h(z)|$ at the same complex
argument, as written on scan~112.  For example, $f=1$ and $h=1+z$
satisfy the coefficient inequality but violate the latter inequality at
$z=-1$.  These corrections repair the proof of the stated proposition.
The uncertain sum/product sign in the transcription of scan~111 is not
used as a premise: the canonical-product hypothesis is stated explicitly.

\paragraph{Finding F-SM-12 (the Fredholm assertions on scans 079--080).}
Let $T$ be a trace-class operator on a complex Hilbert space.  Put
$D_T(z)=\sum_{n\geq0}\alpha_n(T)z^n$, where
$\alpha_n(T)=\operatorname{Tr}(\bigwedge^nT)$ and $\alpha_0(T)=1$.
If $s_j$ are the singular values and $\lambda_j$ the nonzero eigenvalues,
each repeated with its multiplicity (algebraic multiplicity for
$\lambda_j$), then
\[
 |\alpha_n(T)|\leq\alpha_n(|T|),\quad
 |D_T(z)|\leq D_{|T|}(|z|),\quad
 \sum_j|\lambda_j|\leq\|T\|_1,
\]
\[
 D_T(z)=\prod_j(1+z\lambda_j),\qquad
 \operatorname{Tr}T=\sum_j\lambda_j.
\]
In particular $D_T$ has genus zero.  The elementary trace-class ideal,
exterior-power, and compact-operator spectral theorems are the analytic
background here; no arithmetic hypothesis is used.

\begin{proof}
The singular values of $\bigwedge^nT$ are the products
$s_{j_1}\cdots s_{j_n}$ with $j_1<\cdots<j_n$.  Thus
\[
 |\operatorname{Tr}(\bigwedge^nT)|
 \leq\|\bigwedge^nT\|_1
 =\sum_{j_1<\cdots<j_n}s_{j_1}\cdots s_{j_n}
 =\alpha_n(|T|).
\]
The last quantity is at most $\|T\|_1^n/n!$, proving entire convergence
and the two domination inequalities.  Diagonalizing the positive compact
operator $|T|$ gives directly
$D_{|T|}(z)=\prod_j(1+zs_j)$, a genus-zero product.

For any finite collection of nonzero eigenvalues, include their full
generalized eigenspaces and let $E$ be their finite-dimensional invariant
sum.  A unitary triangularization of $T|_E$ gives an orthonormal basis
$e_1,\ldots,e_d$ with diagonal entries the selected eigenvalues.  There is
a diagonal unitary on $E$, extended by zero on $E^\perp$, whose product
with $T$ has trace $\sum_{j=1}^d|\lambda_j|$.  Its operator norm is at
most one.  The trace ideal inequality therefore bounds this sum by
$\|T\|_1$.  Increasing the finite collection proves the summability
assertion.

For completeness, the determinant zero information needed next does not
assume the desired product formula.  The exterior-power definition is
continuous under trace-norm approximation, uniformly for $z$ in compact
sets.  For fixed $n$, this follows by telescoping the tensor product
$A^{\otimes n}-B^{\otimes n}$ and restricting by the bounded
antisymmetrizing projection.  The bounds
$|\alpha_n(A)|\leq\|A\|_1^n/n!$ make the tails uniformly small on
bounded trace-norm sets, proving the asserted locally uniform continuity.
Finite-rank determinant identities consequently imply similarity
invariance and
$\det((1+A)(1+B))=\det(1+A)\det(1+B)$ for trace-class $A,B$.
In particular, if $1+A$ is invertible, its inverse has the form $1+B$
with $B$ trace class, so $\det(1+A)\ne0$.
For a nonzero eigenvalue $\lambda$, the compact spectral theorem gives
the decomposition by its finite-dimensional Riesz subspace and its
invariant closed complement.  Similarity invariance and the direct-sum
identity factor $D_T$ into the ordinary determinant on the Riesz subspace
and the determinant on the complement.  At $z=-1/\lambda$, the latter
factor is nonzero, and the former has zero of order exactly the algebraic
multiplicity of $\lambda$.  At every other point where $1+zT$ is
invertible the determinant is nonzero.  These are all its zeros by the
compact spectral theorem.

Consequently the nonzero zeros of $D_T$ have summable reciprocal moduli.
Its coefficients are dominated by those of the genus-zero function
$D_{|T|}$.  Finding F-SM-11 gives genus zero for $D_T$.  Its zero divisor
and $D_T(0)=1$ identify it as $\prod_j(1+z\lambda_j)$.  Normal convergence
and $\sum_j|\lambda_j|<\infty$ allow differentiation at zero, giving
$\operatorname{Tr}T=\sum_j\lambda_j$.
\end{proof}

This proves the classical Hilbert-space product and trace assertions
written in the deleted passage without using them to prove the
genus-zero lemma.  It should not be described as an RH-dependent solution
of an arithmetic question on those sheets.

\subsection{A fully normalized numerical form of Masullo's main identity}

We now separate a numerical consequence that can be completed from the
adopted Dirichlet theorem from the geometric deconstruction not supplied
in Masullo's Section~12.3.  No kernel-decomposition claim is needed for
the following result.

\paragraph{Finding F-SM-13 (Masullo 12.5.1.1 over $\mathbb Q$, as a numerical distribution).}
Adopt Alp\"oge's Dirichlet zero-location theorem.  Let $\chi$ be a
primitive Dirichlet character of conductor $q$, including the trivial
character of conductor one; write $\chi(-1)=(-1)^a$ with $a\in\{0,1\}$.
Let $h\in C_c^\infty(\mathbb R_{>0})$ satisfy $h(1)=0$, and put
\[
 \widehat h(z)=\int_0^\infty h(t)t^z\,\frac{dt}{t},
 \qquad\delta_\chi=\begin{cases}1,&q=1,\\0,&q>1.\end{cases}
\]
Define the following concrete linear functionals:
\[
 A_a(h)=\frac12\int_0^\infty h(t)t^{1/2}
 \left(\frac1{|1-t|}+\frac{(-1)^a}{1+t}\right)\frac{dt}{t}
 =\int_0^\infty
 \frac{e^{-(a+1/2)x}}{1-e^{-2x}}
 \bigl(h(e^x)+h(e^{-x})\bigr)\,dx,
\]
\[
 B_\chi(h)=\sum_{n\geq1}\frac{\Lambda(n)}{\sqrt n}
 \bigl(\chi(n)h(n)+\overline\chi(n)h(n^{-1})\bigr).
\]
The integrals are ordinary absolutely convergent integrals, and the prime
power sum is finite.  For every $\alpha>0$ one has
\begin{equation}\label{eq:SM-numeric-strip}
 A_a(h)+B_\chi(h)
 =\delta_\chi\bigl(\widehat h(1/2)+\widehat h(-1/2)\bigr)
 -\!\sum_{\substack{\rho:\,L(\rho,\chi)=0,\ 0<\Re\rho<1\\
                         |\Re\rho-1/2|<\alpha}}
       m_\chi(\rho)\widehat h(\rho-1/2).
\end{equation}
The zero sum is absolutely convergent.  Thus this is an equality of
continuous linear functionals on the closed test-function subspace
$\{h\in C_c^\infty(\mathbb R_{>0}):h(1)=0\}$, with its usual inductive
limit topology.

Under the canonical identification
$C_{\mathbb Q}\simeq\widehat{\mathbb Z}^{\times}\times\mathbb R_{>0}$,
take Haar measure $du\,dt/t$, where $du$ has total mass one, and set
\[
 \phi_{\chi,h}(u,t)=\widehat\chi(u)h(t),\qquad
 c_{\chi,s}(u,t)=\widehat\chi(u)^{-1}t^s.
\]
Here $\widehat\chi$ denotes the finite character on
$\widehat{\mathbb Z}^{\times}$ induced by $\chi$.
Here and above $L(s,\chi)$ denotes the finite Dirichlet series.  Masullo's
notation $L(c)$ includes the archimedean factor (Definition~3.2.1.2,
page~123).  With these conventions his finite factor satisfies
$L_f(c_{\chi,s})=L(s,\chi)$, and his completed function is
\[
 L_{\mathrm M}(c_{\chi,s})=
 \pi^{-(s+a)/2}\Gamma((s+a)/2)L(s,\chi).
\]
It has exactly the nontrivial zero divisor appearing above.  Moreover
\[
 M(\phi_{\chi,h})(c_{\chi,s}|\cdot|^{-1/2})
       =\widehat h(s-1/2).
\]
Equation~\eqref{eq:SM-numeric-strip} is exactly the $\chi$-component of
Masullo's Theorem~12.5.1.1, pages~519--520, for $F=\mathbb Q$ and
$I=[-\alpha,\alpha]$.  Finite Fourier decomposition consequently proves
that numerical theorem for every $\phi\in\mathcal D(C_{\mathbb Q})$
with $\phi(C_{\mathbb Q}^1)=0$, in Masullo's own Haar normalization.

\begin{proof}
We give the contour calculation as well as the local normalization checks.
The classical analytic inputs are the continuation and functional equation
of primitive Dirichlet $L$-functions, the order-one completed function,
and its zero count $N_\chi(T)=O_\chi(T\log(2+T))$.  These are background
theorems; this proof does not claim to re-prove the analytic continuation
or the functional equation.  They are the standard inputs behind the
explicit formula (Iwaniec and Kowalski, 2004, Sections~4.6, 5.3, and~5.5).

Let
\[
 \mathcal L_\chi(s)=(q/\pi)^{(s+a)/2}
            \Gamma((s+a)/2)L(s,\chi),\qquad H(s)=\widehat h(s-1/2).
\]
The completed function has precisely the nontrivial zeros, and has simple
poles at $0,1$ exactly when $\delta_\chi=1$.  Its functional equation is
$\mathcal L_\chi(s)=\varepsilon_\chi\mathcal L_{\bar\chi}(1-s)$,
where the nonzero constant $\varepsilon_\chi$ disappears after logarithmic
differentiation.  Integration by parts after $t=e^x$ shows that $H(s)$
decays faster than any inverse power of $|\Im s|$, uniformly on every
fixed vertical strip.  This and the zero count prove absolute convergence
of $\sum_\rho m_\chi(\rho)H(\rho)$.

Integrate $(\mathcal L'_\chi/\mathcal L_\chi)(s)H(s)$ around rectangles
with vertical sides $c=3/2$ and $1-c=-1/2$.  Horizontal sides can be chosen
to contribute zero in the limit.  One explicit justification uses the
Hadamard logarithmic derivative of
$[s(s-1)]^{\delta_\chi}\mathcal L_\chi(s)$: in each interval $[j,j+1]$
choose a height at distance at least $c_0/(j\log(2+j))$ from the ordinates
of all zeros of modulus at most $2j+4$; the zero count leaves such a height
for a sufficiently small fixed $c_0>0$.  Do the same below the real axis.
The nearby terms of the logarithmic derivative are then polynomially
bounded in $j$, and its tail is bounded using
$\sum_{|\rho|>R}|\rho|^{-2}=O_\chi(\log(2+R)/R)$.
The rapid decay of $H$ kills both horizontal integrals.

The residue theorem and the functional equation therefore give
\begin{align*}
 \sum_\rho m_\chi(\rho)H(\rho)-\delta_\chi(H(0)+H(1))
 =\frac1{2\pi i}\int_{(c)}\biggl[
 \frac{\mathcal L'_\chi}{\mathcal L_\chi}(s)H(s)
 +\frac{\mathcal L'_{\bar\chi}}{\mathcal L_{\bar\chi}}(s)H(1-s)
 \biggr]\,ds.
\end{align*}
The constant $\tfrac12\log(q/\pi)$ in each logarithmic derivative
contributes $\log(q/\pi)h(1)=0$, by Mellin inversion.  On this line the
Dirichlet series for $L'/L$ converges absolutely.  Mellin inversion gives
\[
 \frac1{2\pi i}\int_{(c)}H(s)n^{-s}\,ds=n^{-1/2}h(n),\qquad
 \frac1{2\pi i}\int_{(c)}H(1-s)n^{-s}\,ds=n^{-1/2}h(n^{-1}).
\]
Thus the two $L'/L$ terms contribute $-B_\chi(h)$.

For the gamma terms use, for $\Re z>0$,
\[
 \psi(z)=-\gamma+\int_0^\infty
                 \frac{e^{-v}-e^{-zv}}{1-e^{-v}}\,dv.
\]
First truncate the $v$-integral to $[\epsilon,R]$.  All constant terms
pair with $h(1)=0$.  Fubini and Mellin inversion give for the remaining
term
\[
 -\frac12\int_\epsilon^R
 \frac{e^{-(a/2+1/4)v}}{1-e^{-v}}
       \bigl(h(e^{v/2})+h(e^{-v/2})\bigr)\,dv.
\]
The truncated digamma integrals on the line $\Re s=c$ are bounded by
$C(1+|\Im s|)$ uniformly in $\epsilon,R$: near zero use the mean-value
bound on $e^{-v}-e^{-(s+a)v/2}$, and away from zero use exponential
decay.  This justifies passage to the limit against rapidly decreasing
$H(s)+H(1-s)$.  The resulting real integral is absolutely convergent,
because $h(1)=0$ and $h$ is smooth, and equals $-A_a(h)$ after $v=2x$.
The expression for $A_a$ in the statement follows by splitting its
$t$-integral at one and substituting $t=e^x,e^{-x}$; the elementary
identity
\[
 \frac12\left(\frac1{1-r}+\frac{(-1)^a}{1+r}\right)
       =\frac{r^a}{1-r^2},\qquad 0<r<1,
\]
gives exactly its displayed coefficient.  We have proved the full
explicit formula, with the sum over all nontrivial zeros.
The adopted Dirichlet theorem puts every such zero on $\Re\rho=1/2$.
For every $\alpha>0$ its restriction to the indicated open strip therefore
omits no zero, proving~\eqref{eq:SM-numeric-strip}.

We now identify the idelic terms.  For an idele $b=(b_\infty,(b_p)_p)$,
put $r=\operatorname{sgn}(b_\infty)\prod_p p^{v_p(b_p)}\in\mathbb Q^\times$.
Its class has the unique representative
$(t,(u_p)_p)=(b_\infty/r,(b_p/r)_p)$ with $t>0$ and all $u_p$ units;
moreover $t=|b|$.  This proves the stated group identification.
Masullo's Definitions~A.3.2.1--A.3.2.2, page~560, give
$\mu_\infty=dx/|x|$, $\mu_p(\mathbb Z_p^\times)=1$, and
\[
 \mu_1(C_{\mathbb Q}^1)=\frac{2^{r_1}(2\pi)^{r_2}h_{\mathbb Q}R_{\mathbb Q}}
 {\sqrt{|d_{\mathbb Q}|}w_{\mathbb Q}}=\frac{2}{2}=1.
\]
Hence the quotient measure is precisely $du\,dt/t$, without an omitted
volume factor.

For $p\nmid q$, the local idele $b_p=p^kz$ with
$z\in\mathbb Z_p^\times$ has
$\phi_{\chi,h}(b_p)=\chi(p)^{-k}h(p^{-k})$.  The $k=0$ term is zero;
on $k>0$ one has $|1-b_p|_p=1$, and on $k=-m<0$ one has
$|1-b_p|_p=p^m$.  Each shell has multiplicative measure one.  Its local
integral, multiplied by $\log p$, is therefore
\[
 (\log p)\sum_{m\geq1}p^{-m/2}
 \bigl(\chi(p)^m h(p^m)+\overline\chi(p)^m h(p^{-m})\bigr).
\]
For $p\mid q$, primitivity makes the local unit character nontrivial;
unit averaging annihilates each shell with $k\ne0$, while $k=0$ is again
identically zero since $h(1)=0$.  This gives exactly $B_\chi(h)$.
At the real place,
$\phi_{\chi,h}(x)=\operatorname{sgn}(x)^a h(|x|)$; splitting
$\mathbb R^\times$ into its two components and retaining Masullo's
factor $1/2$ gives exactly $A_a(h)$.  The potential singularity at $x=1$
is integrable because $h(1)=0$, so Weil's principal value agrees with the
ordinary integral here.

Orthogonality on the mass-one compact group gives the asserted Mellin
transform, gives zero for all other finite-character components, and gives
$M(\phi_{\chi,h})(|\cdot|^{\pm1/2})=
\delta_\chi\widehat h(\pm1/2)$.  Every Bruhat test function on
$C_{\mathbb Q}$ factors through a finite quotient of its compact factor
and is smooth with compact support on its positive-real factor.  Its
finite Fourier expansion is therefore
$\phi(u,t)=\sum_\chi\widehat\chi(u)h_\chi(t)$, with each character
indexed once by its primitive conductor.  The condition
$\phi(u,1)=0$ for all $u$ is equivalent to $h_\chi(1)=0$ for every
component.  Adding the finitely many component identities proves the
stated assertion for all $\phi$.  Continuity follows on each fixed
compact-support/finite-quotient test space from the local integral bounds,
the finite prime-power sum, and the rapid-decay estimate for the absolutely
convergent zero sum; these estimates are compatible with the inductive
limit topology.
\end{proof}

\paragraph{What is completed and what remains geometric.}
Finding F-SM-13 identifies the numerical distribution in Masullo's
restricted-strip formula for every allowed test function over
$\mathbb Q$, and gives all conductor, parity, Haar, and Mellin
normalizations.  The conductor term vanishes here for the precise reason
$h(1)=0$; it has not been silently discarded in a formula for unrestricted
test functions.  The statement indexes each idelic character by its
primitive conductor; treating an imprimitive Dirichlet series as though it
were that same completed idelic $L$-function would introduce extra deleted
Euler factors.

This numerical identity does not prove Masullo's geometric assertion
\[
 \int_{C_F}k(x,x)\,d\mu(x)
 =\int_{C_F}k_S(x,x)\,d\mu(x)
   +\sum_{p\nmid S}\int_{C_F}k_p(x,x)\,d\mu(x),
\]
nor identify the kernels in the missing Section~12.3, justify an exchange
of the local sum with diagonal restriction/integration, or establish its
Sobolev, nuclearity, and quotient compatibility hypotheses.  Equality of
two already identified numerical distributions does not establish that
a proposed kernel construction realizes them.  It does, however, complete
the exact numerical statement of Theorem~12.5.1.1 for $\mathbb Q$ from
the adopted Alp\"oge theorem and classical analytic inputs.  No extension
to arbitrary number fields or their Hecke characters is made by this
argument.
